% generated by roam.py handbook export -- не редактировать, править заметки
% 20260921124320-теорема_ферма_необходимое_условие_экстремума.org (body ced7b28cf760)

\begin{theorem}[{Теорема Ферма (необходимое условие экстремума)}]\label[theorem]{note:1E03F05F-22CC-4D55-BCF1-EBB5774765F0}
  Пусть $g: (a;b) \rightarrow \R$, $t_0 \in (a;b)$~--- точка локального экстремума $g$, и $g$ дифференцируема в точке $t_0$.
  Тогда $g^{\prime}(t_0) = 0$
\end{theorem}

\begin{remark}
  Важно, что $t_0$~--- \emph{внутренняя} точка: к ней можно подойти с обеих сторон. На конце отрезка теорема неверна: $g(t) = t$ на $[0;1]$ имеет минимум в $t_0 = 0$, но $g^{\prime}(0) = 1$.
\end{remark}

\begin{remark}
  Обратное неверно: у $g(t) = t^3$ производная $g^{\prime}(0) = 0$, но экстремума в нуле нет. Условие $g^{\prime}(t_0) = 0$ только \emph{необходимое}.
\end{remark}

\begin{proof}
  Пусть $t_0$~--- точка локального минимума (для максимума применим то же рассуждение к $-g$).
  Тогда $\exists \delta > 0$: $g(t) \ge g(t_0)$ при $\abs{t - t_0} < \delta$.

  Посмотрим на разностное отношение с двух сторон от $t_0$. Числитель $g(t) - g(t_0) \ge 0$, а знак знаменателя зависит от стороны:
  \[
  t > t_0:\quad \frac{g(t) - g(t_0)}{t - t_0} \ge 0, \qquad t < t_0:\quad \frac{g(t) - g(t_0)}{t - t_0} \le 0
  \]
  Производная в $t_0$ существует, значит оба односторонних предела разностного отношения существуют и равны $g^{\prime}(t_0)$. Нестрогие неравенства сохраняются в пределе:
  \[
  g^{\prime}(t_0) = \lim_{t \rightarrow t_0 + 0} \frac{g(t) - g(t_0)}{t - t_0} \ge 0, \qquad
  g^{\prime}(t_0) = \lim_{t \rightarrow t_0 - 0} \frac{g(t) - g(t_0)}{t - t_0} \le 0
  \]
  Следовательно $g^{\prime}(t_0) = 0$.
\end{proof}
